
qui se vérifie d'%
après 1.10 sur les éléments \(1+Xt,1+Yt\). Pour le premier,
\[
\begin{aligned}
\lambda_u(1+Xt)\circ\lambda_u(1+Yt)
&=(\underline{1}+\{1+Xt\}u)\circ(\underline{1}+\{1+Yt\}u)\\
&=\underline{1}+\{(1+Xt)\circ(1+Yt)\}u\\
&=\underline{1}+\{1+XYt\}u\\
&=\lambda_u((1+Xt)\circ(1+Yt)).
\end{aligned}
\]
Pour le deuxième :
\[
  \lambda_v(\lambda_u(1+Xt))
  =
  \underline{1}+\{\underline{1}+\{1+Xt\}u\}v
  =
  \Ghat^v(\lambda_u)(\lambda_v(1+Xt)).
\]

\item Les exemples géométriques de 2.2 seront traités plus loin.
\end{enumerate}

2.6. On peut généraliser l'exemple 2.5.a) de la façon suivante : Soit \(K\) un anneau commutatif à élément unité, tel que :
\[
\begin{array}{ll}
\text{i)}& K\text{ est sans torsion},\\
\text{ii)}& \forall x\in K,\ \forall n\in\mathbb N,\quad \text{on a}\ \binom{x}{n}\in K
\end{array}
\]
On pose alors, pour tout \(x\in K\), et tout \(n\in\mathbb N\), \(\lambda^n(x)=\binom{x}{n}\), soit :
\begin{equation}\tag{2.6.1}
  \lambda_t(x)=\sum_{n=0}^{\infty}\binom{x}{n}t^n
  =(1+t)^x=\exp(x\Log(1+t)),
\end{equation}
où la dernière identité est considérée dans \((K\otimes_{\mathbb Z}\mathbb Q)[[t]]\). D'après la multiplicativité de l'exponentielle, ceci définit une structure de pré-\(\lambda\)-anneau sur \(K\).

Pour voir que c'est une structure de \(\lambda\)-anneau, il faut prouver les relations 2.4.2, et il suffit de les prouver dans \(K\otimes_{\mathbb Z}\mathbb Q\), d'après 2.6.i). On est alors ramené par fonctorialité au cas \(K=\mathbb Q[X,Y]\), et aux éléments \(X,Y\). Compte tenu de la définition des \(\lambda^i\), les deux membres dans les
